\documentclass[12pt,a4paper]{article}

\usepackage[utf8]{inputenc}
\usepackage[T1]{fontenc}
\usepackage[brazilian]{babel}
\usepackage{lmodern}
\usepackage{amsmath,amssymb}
\usepackage{microtype}

\usepackage{geometry}
\geometry{a4paper, margin=2.5cm}

\usepackage{setspace}
\onehalfspacing
\usepackage{indentfirst}
\setlength{\parindent}{1.25cm}
\setlength{\parskip}{0pt}

\usepackage{booktabs}
\usepackage{array}
\usepackage{siunitx}
\sisetup{output-decimal-marker = {,}, group-digits = false}

\usepackage{caption}
\captionsetup{font=small, labelfont=bf}

\usepackage[hidelinks]{hyperref}

\renewcommand{\contentsname}{Sumário}

\begin{document}

\begin{center}
    {\large Instituto Politécnico do Rio de Janeiro}\\[0.15cm]
    {\large Laboratório de Física 1}\\[0.4cm]
    {\large Experimento 2}\\[0.4cm]
    {\large Aluna: Julia do Amaral Rangel}\\[0.15cm]
    {\large Matrícula: 202310072211}\\[0.6cm]
    {\Large\bfseries Comparação entre os Coeficientes de Atrito Estático e Cinético}\\[0.6cm]
    {\large Nova Friburgo}\\[0.15cm]
    {\large 2026}
\end{center}

\newpage
\tableofcontents
\newpage

\section{Objetivo}

O objetivo do experimento é medir os coeficientes de atrito cinético e
atrito estático de um bloco de madeira que desliza em contato com uma
superfície inclinada. Desta forma, buscamos validar o modelo teórico
comparando os coeficientes obtidos.

\section{Introdução Teórica}

Entendemos o atrito como uma força que surge a partir da interação de um
corpo com uma superfície na qual ele se desloca, seu módulo é
proporcional ao módulo da força normal e sua direção é sempre oposta ao
movimento relativo entre as superfícies (ou à sua tendência). A força de
atrito se faz presente em variadas situações do nosso dia-a-dia, como
caminhar ou escrever em uma folha de papel, e seu estudo se dá desde a
Antiguidade com Aristóteles, porém os estudos mais sistemáticos sobre o
comportamento do atrito são atribuídos a Leonardo da Vinci, que observou
de maneira experimental que a força de atrito era proporcional ao seu
peso e pouco dependia da área de contato. Posteriormente Guillaume
Amontons formalizou as primeiras leis do atrito e, mais tarde
Charles-Augustin de Coulomb trouxe a distinção de atrito estático e
atrito cinético. Desta forma as leis de Amontons-Coulomb permanecem como
um bom modelo para descrever esse fenômeno. Podemos expressar essas
relações a partir da seguinte equação:

\begin{equation}\label{eq:forca_atrito}
F_{at} = \mu\,N
\end{equation}

Onde, \(F\) representa força de atrito, \(N\) representa a força normal
e \(\mu\) representa o coeficiente de atrito estático ou cinético. O
atrito estático se manifesta ao impedir o movimento relativo entre
superfícies, esta força é variável e depende da inclinação da superfície
(\(\theta\)), no experimento estamos procurando a maior força de atrito
possível, aquela que se manifesta na maior inclinação possível que
preserva o estado de inercia do bloco. A partir de decomposições
geométricas podemos relacionar estas entidades a partir da seguinte
expressão:

\begin{equation}\label{eq:mu_estatico}
\mu_e = \tan(\theta_c)
\end{equation}

Onde, \(\mu_{e}\) representa o coeficiente máximo de atrito estático e
\(\theta_{c}\) é o ângulo crítico em que o movimento se inicia
possibilita a determinação do coeficiente de atrito estático. Já no caso
do atrito cinético, consideramos que o corpo está em movimento sobre o
plano inclinado, de forma que a força de atrito possui módulo
aproximadamente constante e atua se opondo ao sentido do movimento.
Aplicando a Segunda Lei de Newton \((F_{r} = m.a)\) na direção paralela
ao plano inclinado, temos:

\begin{equation}\label{eq:segunda_lei}
mg\sin\theta - \mu_c\,(mg\cos\theta) = ma
\end{equation}

\begin{equation}\label{eq:mu_cinetico}
\mu_c = \frac{g\sin\theta - a}{g\cos\theta}
\end{equation}

Como a aceleração medida de forma experimental pode apresentar pequenas
variações, cada valor de aceleração fornece um valor correspondente para
o coeficiente de atrito cinético. Portanto, o coeficiente de atrito
cinético é determinado pela média dos coeficientes obtidos:

\begin{equation}\label{eq:media}
\overline{\mu}_c = \frac{1}{N}\sum_{i=1}^{N} \mu_{c,i}
\end{equation}

Onde \(N\) é o número de medições realizadas. A dispersão dos resultados
é estimada pelo desvio padrão amostral:

\begin{equation}\label{eq:desvio_padrao}
s_{\mu_c} = \sqrt{\frac{1}{N-1}\sum_{i=1}^{N}(\mu_{c,i} - \overline{\mu}_c)^2}
\end{equation}

Que representa a incerteza associada às medidas experimentais do
coeficiente de atrito cinético. Experimentalmente esperamos que o
coeficiente de atrito cinético seja menor que o coeficiente de atrito
estático (\(\mu_e > \mu_c\)).

\section{Materiais Utilizados e Roteiro Experimental}

\subsection{Materiais}

\begin{itemize}
\item Plano Inclinado;
\item Bloco de madeira;
\item Tripé;
\item Câmera do Celular;
\item Software Tracker.
\end{itemize}

\subsection{Execução}

Na primeira etapa do experimento escolhemos começar o estudo pelo atrito
estático, para isso posicionamos o bloco de madeira sobre o plano
inclinado inicialmente na horizontal, e aos poucos aumentamos a
inclinação do plano até que o bloco de madeira superasse o atrito e
começasse a deslizar. Já para estudar o atrito cinético, apenas
mantivemos o ângulo no qual o bloco havia começado a deslizar na etapa
anterior, e assim deixamos com que o corpo deslizasse até o final do
plano.

As duas fases do experimento foram gravadas com a câmera do celular e
posteriormente analisadas no software Tracker.

\section{Resultados}

\begin{table}[htbp]
\centering
\caption{Posição $x$ e aceleração $a$ do bloco de madeira em função do tempo $t$.}
\label{tab:dados}
\begin{tabular}{S[table-format=-1.3] S[table-format=1.3] S[table-format=1.3]}
\toprule
{$x$ (m)} & {$a$ (m/s$^2$)} & {$t$ (s)} \\
\midrule
-0.127 & {}      & 0     \\
-0.135 & {}      & 0.033 \\
-0.147 & 1.579   & 0.067 \\
-0.16  & 0.762   & 0.1   \\
-0.174 & 0.173   & 0.133 \\
-0.188 & 0.16    & 0.167 \\
-0.202 & 0.294   & 0.2   \\
-0.215 & 0.614   & 0.233 \\
-0.23  & 1.111   & 0.267 \\
-0.246 & 0.838   & 0.3   \\
-0.264 & 1.11    & 0.333 \\
-0.281 & 1.746   & 0.367 \\
-0.301 & 2.723   & 0.4   \\
-0.324 & 2.16    & 0.433 \\
-0.35  & 1.896   & 0.467 \\
-0.377 & 2.181   & 0.5   \\
-0.407 & 2.627   & 0.533 \\
-0.44  & 2.302   & 0.567 \\
-0.475 & {}      & 0.6   \\
-0.513 & {}      & 0.633 \\
\bottomrule
\end{tabular}
\end{table}

Para determinar o coeficiente de atrito estático, utilizamos o ângulo
crítico \((\theta = 15^\circ)\) medido no laboratório na
Eq.~\eqref{eq:mu_estatico} e chegamos no resultado
\(\mu_{e} = 0,267\). Já para o cálculo do coeficiente de atrito
cinético, usamos o ângulo \((\theta = 34,8^\circ)\), consideramos a
gravidade como \((g = 9,8\ \text{m/s}^2)\), mas não podemos substituir
direto na Eq.~\eqref{eq:mu_cinetico} pois a aceleração apresenta
pequenas variações ao longo do deslocamento, então calculamos o
coeficiente para cada aceleração na Eq.~\eqref{eq:media} e o desvio na
Eq.~\eqref{eq:desvio_padrao}, dessa forma temos que o coeficiente de
atrito cinético é \(\mu_{c} = 0,522 \pm 0,115\).

Para deixar clara a escolha dos ângulos, cabe destacar que cada ângulo
cumpre uma função distinta no experimento. O ângulo de \(15^\circ\) é o
ângulo crítico, isto é, a menor inclinação necessária para que o bloco
vença o atrito estático e inicie o movimento; por isso, é ele que
determina o coeficiente de atrito estático pela Eq.~\eqref{eq:mu_estatico}.
O ângulo de \(34,8^\circ\), por sua vez, é a inclinação utilizada na
etapa do atrito cinético, na qual o bloco percorre o plano com uma
aceleração mensurável, permitindo o cálculo de \(\mu_c\) pelas
Eqs.~\eqref{eq:mu_cinetico} e \eqref{eq:media}. Como a comparação entre
\(\mu_e\) e \(\mu_c\) só é fisicamente consistente quando ambos são
avaliados sob as mesmas condições de inclinação, os dois coeficientes
foram referidos ao mesmo ângulo de \(34,8^\circ\).

Neste experimento temos como objetivo comparar o coeficiente de atrito
estático com o coeficiente de atrito cinético, mas para isso precisamos
comparar para um mesmo ângulo, vamos então comparar os coeficientes no
ângulo \((\theta = 34,8^\circ)\). Como já calculamos o coeficiente de
atrito cinético para esse ângulo \(\mu_{c} = 0,522 \pm 0,115\),
precisamos apenas substituir esse mesmo ângulo na
Eq.~\eqref{eq:mu_estatico}, assim temos \(\mu_{e} = 0,695\). Logo, o
resultado obtido atende o que se espera pela teoria
(\(\mu_e > \mu_c\)).

\section{Conclusão}

O experimento permitiu determinar os coeficientes de atrito estático e
cinético entre o bloco de madeira e o plano inclinado. A partir do
ângulo crítico, obteve-se o coeficiente de atrito estático, enquanto o
coeficiente de atrito cinético foi calculado com base nas acelerações
medidas pelo software Tracker durante o deslizamento do bloco. Os
resultados encontrados foram \(\mu_{e} = 0,695\) e
\(\mu_{c} = 0,522 \pm 0,115\), verificando concordância com o modelo
teórico \(\mu_e > \mu_c\). As pequenas variações observadas nos valores
de aceleração ao longo do movimento foram consideradas no cálculo do
coeficiente de atrito cinético por meio da média e do desvio padrão.
Dessa forma, concluímos que o experimento atingiu seu objetivo,
reproduzindo qualitativamente o comportamento esperado para o atrito
entre superfícies sólidas.

\end{document}
